\section{Свойство линейности для собственных векторов. Собственное подпространство. Критерий диагональности матрицы оператора. Примеры}

\begin{definition}
Пусть $\lambda$ --- собственное число линейного оператора $\varphi \in L(V, V)$. Множество всех собственных векторов, соответствующих $\lambda$, вместе с нулевым вектором, называется \textbf{собственным подпространством} оператора $\varphi$, соответствующим $\lambda$, и обозначается $L_\lambda$:
\[
L_\lambda = \{ x \in V \mid \varphi(x) = \lambda x \} = \ker(\varphi - \lambda I),
\]
где $I$ --- тождественный оператор.
\end{definition}

\begin{remark}
Поскольку ядро любого линейного оператора является подпространством, $L_\lambda$ действительно является подпространством пространства $V$.
\end{remark}

\begin{theorem}[О свойстве линейности собственных векторов]
Если $x_1, x_2, \dots, x_k$ --- собственные векторы оператора $\varphi$, соответствующие \textbf{одному и тому же} собственному числу $\lambda_0$, то любая их ненулевая линейная комбинация $y = c_1 x_1 + c_2 x_2 + \dots + c_k x_k$ ($c_i \in F$, не все равны нулю) также является собственным вектором оператора $\varphi$, соответствующим $\lambda_0$.
\end{theorem}

\begin{proof}
Применим оператор $\varphi$ к вектору $y$ и воспользуемся его линейностью:
\[
\varphi(y) = \varphi(c_1 x_1 + \dots + c_k x_k) = c_1 \varphi(x_1) + \dots + c_k \varphi(x_k).
\]
Так как все $x_i$ --- собственные векторы для $\lambda_0$, то $\varphi(x_i) = \lambda_0 x_i$. Подставляем это в равенство:
\[
\varphi(y) = c_1 (\lambda_0 x_1) + \dots + c_k (\lambda_0 x_k) = \lambda_0 (c_1 x_1 + \dots + c_k x_k) = \lambda_0 y.
\]
Поскольку по условию $y \neq 0$, вектор $y$ удовлетворяет определению собственного вектора с собственным числом $\lambda_0$.
\hfill$\blacksquare$
\end{proof}

\begin{definition}
Квадратная матрица $D$ порядка $n$ называется \textbf{диагональной}, если все её элементы вне главной диагонали равны нулю. Она имеет вид:
\[
D = \operatorname{diag}(\lambda_1, \lambda_2, \dots, \lambda_n) = 
\begin{pmatrix}
\lambda_1 & 0 & \dots & 0 \\
0 & \lambda_2 & \dots & 0 \\
\vdots & \vdots & \ddots & \vdots \\
0 & 0 & \dots & \lambda_n
\end{pmatrix}.
\]
Оператор $\varphi$ называется \textbf{диагонализируемым}, если существует базис пространства $V$, в котором его матрица $A_\varphi$ является диагональной.
\end{definition}

\begin{theorem}[Критерий диагонализируемости оператора]
Линейный оператор $\varphi$ диагонализируем тогда и только тогда, когда в пространстве $V$ существует базис, состоящий целиком из собственных векторов этого оператора. При этом на главной диагонали матрицы оператора в этом базисе стоят соответствующие собственные числа.
\end{theorem}

\begin{proof}
\begin{itemize}
    \item \textbf{Необходимость:} Пусть в некотором базисе $E$ матрица оператора диагональна: $A_\varphi = \operatorname{diag}(\lambda_1, \dots, \lambda_n)$. Это означает, что $\varphi(e_i) = \lambda_i e_i$ для каждого базисного вектора $e_i$. Следовательно, все векторы базиса $E$ являются собственными.
    \item \textbf{Достаточность:} Пусть $E = \{e_1, \dots, e_n\}$ --- базис из собственных векторов, и $\varphi(e_i) = \lambda_i e_i$. Тогда по определению матрицы оператора, её $i$-й столбец состоит из координат вектора $\lambda_i e_i$ в базисе $E$, то есть это столбец, у которого на $i$-м месте стоит $\lambda_i$, а остальные элементы равны нулю. Следовательно, $A_\varphi$ --- диагональная матрица.
\end{itemize}
\hfill$\blacksquare$
\end{proof}

\begin{example}
Рассмотрим оператор в $\mathbb{R}^3$ с матрицей в стандартном базисе:
\[
A = \begin{pmatrix}
2 & 1 & 0 \\
0 & 2 & 0 \\
0 & 0 & 2
\end{pmatrix}.
\]
Характеристическое уравнение: $\det(A - \lambda E) = (2-\lambda)^3 = 0$. Единственное собственное число $\lambda = 2$ (алгебраическая кратность 3).
Найдем собственные векторы, решая систему $(A - 2E)X = 0$:
\[
\begin{pmatrix}
0 & 1 & 0 \\
0 & 0 & 0 \\
0 & 0 & 0
\end{pmatrix}
\begin{pmatrix} x_1 \\ x_2 \\ x_3 \end{pmatrix} = \begin{pmatrix} 0 \\ 0 \\ 0 \end{pmatrix} \quad \Rightarrow \quad x_2 = 0.
\]
Переменные $x_1$ и $x_3$ --- свободные. Фундаментальная система решений состоит из двух векторов: $v_1 = (1, 0, 0)^T$ и $v_2 = (0, 0, 1)^T$. 
Так как геометрическая кратность (2) меньше алгебраической (3), базиса из собственных векторов не существует, и оператор не диагонализируем.
\end{example}
